% ============================================================================
% 서식 설정 — main.tex 에서 \input 한다.
% ============================================================================
\usepackage{iftex}
\ifXeTeX\else\ifLuaTeX\else
  \errmessage{이 문서는 한글 조판 때문에 XeLaTeX(권장) 또는 LuaLaTeX로 컴파일해야 합니다.}
\fi\fi

% ── 한글 ────────────────────────────────────────────────────────────────
\usepackage{kotex}
\usepackage{fontspec}
% 설치된 글꼴을 순서대로 찾는다: Noto CJK(Overleaf·대부분 리눅스) → 나눔 → 은(TeX Live 기본)
\IfFontExistsTF{Noto Serif CJK KR}{\setmainhangulfont{Noto Serif CJK KR}[AutoFakeSlant=0.15]}{%
  \IfFontExistsTF{NanumMyeongjo}{\setmainhangulfont{NanumMyeongjo}[AutoFakeSlant=0.15]}{\setmainhangulfont{UnBatang}[AutoFakeSlant=0.15]}}
\IfFontExistsTF{Noto Sans CJK KR}{\setsanshangulfont{Noto Sans CJK KR}}{%
  \IfFontExistsTF{NanumGothic}{\setsanshangulfont{NanumGothic}}{\setsanshangulfont{UnDotum}}}
\IfFontExistsTF{Noto Sans Mono CJK KR}{\setmonohangulfont{Noto Sans Mono CJK KR}}{%
  \IfFontExistsTF{NanumGothicCoding}{\setmonohangulfont{NanumGothicCoding}}{\setmonohangulfont{UnDotum}}}

% ── 쪽 ──────────────────────────────────────────────────────────────────
\usepackage[a4paper,top=30mm,bottom=30mm,left=30mm,right=25mm,headheight=14pt]{geometry}
\usepackage{setspace}
\setstretch{1.55}
\setlength{\parindent}{1em}
\setlength{\parskip}{0.25em}
\usepackage{fancyhdr}
\pagestyle{plain}

% ── 수식·표·그림 ────────────────────────────────────────────────────────
\usepackage{amsmath,amssymb,mathtools,bm}
\usepackage{graphicx}
\graphicspath{{figures/}}
\usepackage{booktabs,multirow,array}
\usepackage{adjustbox}
\usepackage{float}
\usepackage[labelsep=period,font=small,labelfont=bf,skip=6pt]{caption}
\usepackage{subcaption}
\usepackage[section]{placeins}   % 절 경계를 넘어 그림·표가 떠다니지 않게
\usepackage{algorithm}
\usepackage{algpseudocode}
\usepackage{enumitem}
\setlist{itemsep=0.15em,topsep=0.4em,leftmargin=1.6em}
\usepackage{xcolor}
\usepackage{url}

% ── 참고문헌·링크 ──────────────────────────────────────────────────────
\usepackage[numbers,sort&compress]{natbib}
\setlength{\bibsep}{0.3em}
\usepackage[unicode,hidelinks,bookmarksnumbered]{hyperref}

% ── 장·절 제목 ─────────────────────────────────────────────────────────
\usepackage{titlesec}
\titleformat{\chapter}[display]{\centering\bfseries}{\Large 제\,\thechapter\,장}{0.5em}{\LARGE}
\titlespacing*{\chapter}{0pt}{0pt}{2.2em}
\titleformat{\section}{\bfseries\large}{\thesection}{0.8em}{}
\titleformat{\subsection}{\bfseries\normalsize}{\thesubsection}{0.8em}{}
\titlespacing*{\section}{0pt}{1.4em}{0.6em}
\titlespacing*{\subsection}{0pt}{1.0em}{0.4em}
\setcounter{secnumdepth}{2}
\setcounter{tocdepth}{2}

% ── 한국어 명칭 (kotex 가 문서 시작 때 덮어쓰지 않도록 \AtBeginDocument 에서 설정) ──
\AtBeginDocument{%
  \renewcommand{\contentsname}{목\quad 차}%
  \renewcommand{\listfigurename}{그림 목차}%
  \renewcommand{\listtablename}{표 목차}%
  \renewcommand{\bibname}{참고문헌}%
  \renewcommand{\figurename}{그림}%
  \renewcommand{\tablename}{표}%
  \renewcommand{\appendixname}{부록}%
}
\floatname{algorithm}{알고리즘}
\algrenewcommand\algorithmicrequire{\textbf{입력:}}
\algrenewcommand\algorithmicensure{\textbf{출력:}}
% 부록 장 제목은 '부록 A'
\usepackage{etoolbox}
\apptocmd{\appendix}{%
  \titleformat{\chapter}[display]{\centering\bfseries}{\Large 부록\,\thechapter}{0.5em}{\LARGE}}{}{}

% ── 본문용 매크로 ──────────────────────────────────────────────────────
\newcommand{\fillin}[1]{\textcolor{red}{#1}}            % 저자가 채울 자리 (최종본에서 빨간색 확인)
\newcommand{\TODO}[1]{\textcolor{red}{[확인 필요: #1]}}  % 저자가 확인해야 할 서술
\DeclareRobustCommand{\figref}[1]{그림~\ref{#1}}
\DeclareRobustCommand{\tabref}[1]{표~\ref{#1}}
\DeclareRobustCommand{\secref}[1]{\ref{#1}절}
\DeclareRobustCommand{\chapref}[1]{제\ref{#1}장}
\DeclareRobustCommand{\eqnref}[1]{식~(\ref{#1})}
\DeclareRobustCommand{\algref}[1]{알고리즘~\ref{#1}}
\DeclareRobustCommand{\appref}[1]{부록~\ref{#1}}
\newcommand{\BRM}{SegmentBRM}
\newcommand{\mAP}{mAP$_{50:95}$}
\newcommand{\tn}{TN3K}
\newcommand{\dd}{DDTI}
\newcommand{\todir}{$\rightarrow$}
\DeclareMathOperator*{\argmax}{arg\,max}
\DeclareMathOperator{\BCE}{BCE}
\DeclareMathOperator{\dil}{dilate}
\DeclareMathOperator{\ero}{erode}
